\section{Vulnerability and Defense Context}
\label{sec:vuln}

\subsection{Root Cause: Data-Dependent Power Consumption}

CMOS logic gates consume dynamic power proportional to switching activity.
When a register or bus transitions between logic states, the current drawn
from the supply correlates with the Hamming distance (HD) between the previous
and new value---and when the previous state can be assumed constant (e.g.,
after a reset), this reduces to a correlation with the Hamming weight (HW) of
the new value.

In AES-128, the \textbf{SubBytes} (S-Box) operation in round 1 maps each byte
$p_b$ of plaintext to $S(p_b \oplus k_b)$, where $k_b$ is the corresponding
secret key byte.
Because the S-Box input is a function of a \emph{known} plaintext byte XORed
with an \emph{unknown} key byte, an adversary can:

\begin{enumerate}
  \item Hypothesise all 256 possible values of $k_b$.
  \item Compute the predicted intermediate value $v_{d,j} = S(p_{d,b} \oplus j)$
        for each trace $d$ and hypothesis $j$.
  \item Derive the predicted power $h_{d,j} = \mathrm{HW}(v_{d,j})$.
  \item Use a statistical distinguisher to find which $j^*$ explains the
        observed power the best---that $j^*$ is the key byte.
\end{enumerate}

Attacking the \emph{last round} (AddRoundKey$_{10}$) is symmetrically
possible using ciphertexts and the inverse S-Box:
$v_{d,j} = S^{-1}(c_{d,b} \oplus j)$,
where $j$ now hypothesises bytes of the round-10 key.
The original key is then recovered via backward key schedule.
This is the approach used on the reference dataset in this project (see
Section~\ref{sec:method:attack}).

\subsection{Leakage Models}

Two leakage models are evaluated:

\begin{description}
  \item[Hamming Weight (HW).]
    $h_{d,j} = \mathrm{HW}(v_{d,j})$, i.e., the number of set bits in the
    intermediate value.
    This assumes the power consumed per operation is proportional to the
    number of 1-bits processed.

  \item[Hamming Distance (HD).]
    $h_{d,j} = \mathrm{HW}(v_{d,j} \oplus c)$ for a constant $c$, modelling
    transitions from a previous state value $c$.
    In the CPA implementation, $c$ (called \texttt{const} in the code) is
    swept over $\{0, \ldots, 255\}$ and the first value that allows full key
    recovery is reported.
\end{description}

\subsection{Standard Countermeasures}

The primary defence against power side-channel attacks is \textbf{masking}:
splitting the sensitive intermediate variable $v = S(p \oplus k)$ into $n$
shares $v_1, \ldots, v_n$ such that $v = v_1 \oplus \cdots \oplus v_n$, where
each share is individually statistically independent of the key.
A first-order CPA attack fails against a first-order masked implementation
because no single power trace sample correlates with $k$; a second-order
attack combining two samples is required, exponentially increasing trace
complexity.

Evaluating masked implementations is proposed as future work
(Section~\ref{sec:conclusion}).
The current experiments target an \emph{unprotected} implementation to
establish baseline attack efficiency.
